%% ma: L11-B17-0065
%% lop: 11
%% chuong: 5
%% bai: 17
%% dang: 11.17.5
%% loai: DS
%% muc: [TH, TH, VD, VD]
%% dap_an: "ĐĐSS"
%% nguon: "tai-lieu-hoc-tap-mon-toan-lop-11-theo-chuong-trinh-giao-duc-pho-thong-2018.pdf (Lớp Toán Thầy Duy, Phạm Lê Duy, THPT Chu Văn An), Chương 3 Bài 3, Phần 2, Câu 5 (tr.417)"
%% kien_thuc: "Lớp 11 Bài 17 (hàm số liên tục)"
%% trang_thai: cho-duyet
%% ly_do: "Dạng mới đề xuất (11.17.x), chờ giáo viên duyệt dạng; tài liệu gốc không có lời giải mẫu, đáp án do kho tự giải và kiểm bằng sympy."
\begin{ex}
Cho $a,b$ là hai số thực sao cho hàm số $f(x)=\begin{cases}\dfrac{x^2+ax+b}{x-1}&\text{khi }x\ne1\\2ax-1&\text{khi }x=1\end{cases}$ liên tục trên $\mathbb R$. Khi đó
\choiceTF
{\True $f(1)=2a-1$.}
{\True $a>0$.}
{$b>0$.}
{$a-b=6$.}
\loigiai{\sloppy Giới hạn $\lim\limits_{x\to1}\dfrac{x^2+ax+b}{x-1}$ hữu hạn nên $1+a+b=0$, khi đó $x^2+ax+b=(x-1)(x+1+a)$ và giới hạn bằng $2+a$. Liên tục: $2+a=2a-1\Rightarrow a=3$, $b=-4$.
a) $f(1)=2a-1$. Đúng. b) $a=3>0$. Đúng. c) $b=-4<0$. Sai. d) $a-b=7\ne6$. Sai.}
\end{ex}
